\section{Data Sources and Technical Environment}
\label{sec:data-sources}

This study integrates existing soil and ancillary
spatial datasets with Sentinel-2 satellite imagery
to examine soil-map coverage, vegetation
distribution, and temporal vegetation dynamics
within the study area which were obtained
from the Sentinel-2 Level-2A product collection
through the Copernicus Data Space Ecosystem (CDSE) \mkbibbrackets{\cite{drusch2012sentinel2}}. The workflow combines
vector-data preparation and spatial analysis
with satellite image processing and NDVI-based
assessment.

The principal datasets, their spatial
characteristics, and their roles in the analysis
are summarized in
\Cref{tab:data-sources}.

\begin{table}[H]
	\centering
	\caption{Summary of spatial and satellite data
		sources used in the GeoSoil-Cuba workflow.}
	\label{tab:data-sources}
	\small
	\begin{tabularx}{\textwidth}{
			@{}
			p{2.6cm}
			p{2.5cm}
			p{2.8cm}
			X
			@{}
		}
		\toprule
		\textbf{Dataset} &
		\textbf{Source or format} &
		\textbf{Spatial reference} &
		\textbf{Role in the study} \\
		\midrule
		
		Soil map &
		Faculty of Geography, University of Havana;
		MapInfo TAB &
		NAD 27 Cuba datum; Lambert Conformal Conic;
		metres &
		Delineation and analysis of mapped soil units
		within the study area. Original scale:
		1:25,000. \\
		
		Area of Interest (AOI) &
		KML boundary &
		WGS 84 geographic coordinates
		(EPSG:4326) &
		Defines the study-area boundary and spatial
		extent of the analysis. \\
		
		Water bodies &
		Faculty of Geography, University of Havana;
		MapInfo TAB &
		As supplied; harmonized for spatial analysis &
		Identification of water surfaces and
		interpretation of apparent gaps in soil-map
		coverage. \\
		
		Vegetation and land cover &
		Faculty of Geography, University of Havana;
		vector dataset &
		As supplied; harmonized for spatial analysis &
		Integration of soil units with vegetation
		and land-cover attributes. \\
		
		Sentinel-2 imagery &
		Copernicus Data Space Ecosystem;
		Level-2A products &
		Tiles T17QLF and T17QMF;
		bands B04 and B08 at 10 m resolution;
		SCL at 20 m &
		Satellite-based vegetation analysis,
		including NDVI calculation and temporal
		assessment. Search period: November 2025
		to June 2026. \\
		
		\bottomrule
	\end{tabularx}
\end{table}


\subsection{Soil and Ancillary Spatial Data}
\label{subsec:soil-data}

The primary soil dataset consists of a
1:25,000-scale soil map of Mayabeque Province,
Cuba, distributed in MapInfo format. The dataset
contains polygon geometries representing mapped
soil units and associated attributes describing
soil types and subtypes according to the
classification system used in the source database.

The soil map uses the NAD27 / Cuba Norte
projected coordinate reference system
(EPSG:3795). Its original projection definition
was examined and standardized before spatial
processing.

The Area of Interest (AOI) was defined using a
polygon provided in Keyhole Markup Language
(KML) format, originally referenced to WGS 84
(EPSG:4326). The AOI was transformed to the
soil map's projected coordinate reference system
for the relevant spatial operations.

An additional water-body dataset, also supplied
in MapInfo format, was incorporated to identify
the Mampostón reservoir and distinguish mapped
water surfaces from areas lacking soil-map
coverage.

The original spatial datasets were preserved
unchanged. Derived geometries and attributes
were stored separately to maintain traceability
throughout the analytical workflow.

The original production dates and authorship
of the soil and ancillary datasets have not
been independently verified. Their spatial
and thematic interpretation must therefore
account for the characteristics and limitations
of the supplied source data.


\subsection{Vegetation and Land-Cover Data}
\label{subsec:land-cover-data}

Vegetation and land-cover information was obtained
from an existing vector dataset containing
categorical descriptions of vegetation and
land-cover units.

The dataset includes the attributes
\texttt{LCNAME}, \texttt{NEWVEG}, and
\texttt{TNC\_CLASS}, which provide complementary
classification information.

These attributes were examined and incorporated
into the spatial integration workflow to
characterize vegetation and land-cover
categories associated with the mapped soil units.

The vegetation dataset was processed independently
from the Sentinel-2 imagery. Its categorical
classifications therefore represent information
from an existing spatial dataset rather than
land-cover classes derived from the satellite
observations analyzed in this study.

The production date and original classification
methodology of the vegetation dataset have not
been independently verified. Consequently,
its temporal correspondence with the
2025--2026 satellite observations cannot be
assumed.

The soil, water-body, and vegetation-cover
datasets were provided courtesy of the Faculty
of Geography, University of Havana, Cuba
(Facultad de Geografía, Universidad de La Habana),
with permission for research, analysis, and
public redistribution. The institution is
acknowledged as the data provider; the original
production dates and authorship of these datasets
have not been independently verified.


\subsection{Sentinel-2 Satellite Imagery}
\label{subsec:sentinel-data}

Multispectral satellite observations were obtained
from the Sentinel-2 Level-2A product collection
through the Copernicus Data Space Ecosystem
(CDSE).

Satellite products were discovered using the
SpatioTemporal Asset Catalog (STAC) interface,
which supports searches based on geographic
extent, acquisition date, and product metadata.

The catalogue search covered the period from
November 2025 through June 2026. Products
associated with Sentinel-2 tiles T17QLF and
T17QMF were considered for the study area.

The analysis focused on the red (B04) and
near-infrared (B08) spectral bands, both provided
at a spatial resolution of 10 m. The Scene
Classification Layer (SCL), available at a
spatial resolution of 20 m, was used to support
quality screening and the assessment of
potential cloud and shadow contamination.

The initial satellite inventory contained
226 observations. Following image processing,
32 NDVI rasters were prepared for the selected
agricultural analysis region. A subsequent
regional cloud-screening procedure retained
15 acquisitions satisfying the criterion of
less than 10\% cloud coverage within that region.

This regional criterion reduced the influence
of cloud contamination but did not establish
that all pixels in the retained observations
were free of clouds, haze, or shadows.

Detailed acquisition selection, radiometric
correction, quality screening, and NDVI
calculation procedures are described in
\Cref{sec:methodology}.


\subsection{Software and Technical Environment}
\label{subsec:software}

GeoSoil-Cuba was developed using an open-source
Python geospatial environment managed through
Conda, with packages obtained primarily from
the \texttt{conda-forge} channel.

The documented environment uses Python 3.12.14
and GDAL 3.13.3.

Vector-data processing was performed using
GeoPandas 1.1.4 and Shapely 2.1.2, while
coordinate reference system transformations
were supported by PyProj 3.8.0.

Raster processing was supported by Rasterio
1.5.1, including the handling of georeferenced
satellite imagery and derived NDVI products.

NumPy 2.5.3 and Pandas 3.0.5 provided numerical
and tabular data processing capabilities,
respectively. Matplotlib 3.11.1 was used to
generate analytical figures and visualizations.

Sentinel-2 catalogue discovery was supported
by PySTAC Client 0.9.0, while Requests 2.34.2
provided HTTP communication functionality.

The analytical workflow was organized into
Jupyter notebooks, with JupyterLab 4.6.3
included in the documented environment.
QGIS 3.44.14-Solothurn was also used for the visual selection
and digitization of the three potato-field
sampling polygons.

The Conda environment specification and project
source code provide a record of the principal
software dependencies used in the analysis.
The software versions reported here correspond
to the documented environment and should be
verified against the final environment
specification before publication.


\subsection{Data Characteristics and Limitations}
\label{subsec:data-limitations}

The integration of historical soil cartography,
existing vegetation classifications, and
satellite observations introduces differences
in spatial representation, classification
structure, and temporal coverage.

The soil dataset represents mapped soil units
rather than direct measurements of contemporary
soil conditions. Its interpretation is therefore
constrained by the characteristics of the source
cartographic dataset, including its original
scale, soil classification system, and spatial
completeness. The dataset uses a Lambert
Conformal Conic projection referenced to the
NAD 27 Cuba datum, with coordinates expressed
in metres. The original map scale and
classification metadata should be considered
when interpreting the spatial detail and
thematic meaning of the mapped soil units.

The vegetation dataset provides categorical
land-cover information, whereas Sentinel-2
observations describe spectral conditions
at specific acquisition dates. Because the
original production date of the vegetation
dataset has not been verified, these sources
cannot be assumed to represent the same period.

Satellite-derived vegetation indices may also
be affected by atmospheric conditions, residual
cloud or shadow contamination, land-cover
heterogeneity, radiometric processing choices,
and the spatial resolution of the imagery.

In addition, radiometric parameters were not
independently verified for every processed
Sentinel-2 product. The use of common parameters
for products with unavailable metadata is
therefore a methodological assumption that
should be considered when interpreting the
derived NDVI values.

These characteristics were considered when
preparing the datasets and interpreting the
spatial coverage results and temporal NDVI
patterns. The combined workflow supports
geospatial integration and exploratory
vegetation monitoring but does not, by itself,
establish causal relationships between soil
properties, vegetation conditions, and
agricultural management.